% Options for packages loaded elsewhere
% Options for packages loaded elsewhere
\PassOptionsToPackage{unicode}{hyperref}
\PassOptionsToPackage{hyphens}{url}
%
\documentclass[
  ignorenonframetext,
]{beamer}
\newif\ifbibliography
\usepackage{pgfpages}
\setbeamertemplate{caption}[numbered]
\setbeamertemplate{caption label separator}{: }
\setbeamercolor{caption name}{fg=normal text.fg}
\beamertemplatenavigationsymbolsempty
% remove section numbering
\setbeamertemplate{part page}{
  \centering
  \begin{beamercolorbox}[sep=16pt,center]{part title}
    \usebeamerfont{part title}\insertpart\par
  \end{beamercolorbox}
}
\setbeamertemplate{section page}{
  \centering
  \begin{beamercolorbox}[sep=12pt,center]{section title}
    \usebeamerfont{section title}\insertsection\par
  \end{beamercolorbox}
}
\setbeamertemplate{subsection page}{
  \centering
  \begin{beamercolorbox}[sep=8pt,center]{subsection title}
    \usebeamerfont{subsection title}\insertsubsection\par
  \end{beamercolorbox}
}
% Prevent slide breaks in the middle of a paragraph
\widowpenalties 1 10000
\raggedbottom
\AtBeginPart{
  \frame{\partpage}
}
\AtBeginSection{
  \ifbibliography
  \else
    \frame{\sectionpage}
  \fi
}
\AtBeginSubsection{
  \frame{\subsectionpage}
}
\usepackage{iftex}
\ifPDFTeX
  \usepackage[T1]{fontenc}
  \usepackage[utf8]{inputenc}
  \usepackage{textcomp} % provide euro and other symbols
\else % if luatex or xetex
  \usepackage{unicode-math} % this also loads fontspec
  \defaultfontfeatures{Scale=MatchLowercase}
  \defaultfontfeatures[\rmfamily]{Ligatures=TeX,Scale=1}
\fi
\usepackage{lmodern}

\ifPDFTeX\else
  % xetex/luatex font selection
\fi
% Use upquote if available, for straight quotes in verbatim environments
\IfFileExists{upquote.sty}{\usepackage{upquote}}{}
\IfFileExists{microtype.sty}{% use microtype if available
  \usepackage[]{microtype}
  \UseMicrotypeSet[protrusion]{basicmath} % disable protrusion for tt fonts
}{}
\makeatletter
\@ifundefined{KOMAClassName}{% if non-KOMA class
  \IfFileExists{parskip.sty}{%
    \usepackage{parskip}
  }{% else
    \setlength{\parindent}{0pt}
    \setlength{\parskip}{6pt plus 2pt minus 1pt}}
}{% if KOMA class
  \KOMAoptions{parskip=half}}
\makeatother


\usepackage{longtable,booktabs,array}
\newcounter{none} % for unnumbered tables
\usepackage{calc} % for calculating minipage widths
\usepackage{caption}
% Make caption package work with longtable
\makeatletter
\def\fnum@table{\tablename~\thetable}
\makeatother
\usepackage{graphicx}
\makeatletter
\newsavebox\pandoc@box
\newcommand*\pandocbounded[1]{% scales image to fit in text height/width
  \sbox\pandoc@box{#1}%
  \Gscale@div\@tempa{\textheight}{\dimexpr\ht\pandoc@box+\dp\pandoc@box\relax}%
  \Gscale@div\@tempb{\linewidth}{\wd\pandoc@box}%
  \ifdim\@tempb\p@<\@tempa\p@\let\@tempa\@tempb\fi% select the smaller of both
  \ifdim\@tempa\p@<\p@\scalebox{\@tempa}{\usebox\pandoc@box}%
  \else\usebox{\pandoc@box}%
  \fi%
}
% Set default figure placement to htbp
\def\fps@figure{htbp}
\makeatother


% definitions for citeproc citations
\NewDocumentCommand\citeproctext{}{}
\NewDocumentCommand\citeproc{mm}{%
  \begingroup\def\citeproctext{#2}\cite{#1}\endgroup}
\makeatletter
 % allow citations to break across lines
 \let\@cite@ofmt\@firstofone
 % avoid brackets around text for \cite:
 \def\@biblabel#1{}
 \def\@cite#1#2{{#1\if@tempswa , #2\fi}}
\makeatother
\newlength{\cslhangindent}
\setlength{\cslhangindent}{1.5em}
\newlength{\csllabelwidth}
\setlength{\csllabelwidth}{3em}
\newenvironment{CSLReferences}[2] % #1 hanging-indent, #2 entry-spacing
 {\begin{list}{}{%
  \setlength{\itemindent}{0pt}
  \setlength{\leftmargin}{0pt}
  \setlength{\parsep}{0pt}
  % turn on hanging indent if param 1 is 1
  \ifodd #1
   \setlength{\leftmargin}{\cslhangindent}
   \setlength{\itemindent}{-1\cslhangindent}
  \fi
  % set entry spacing
  \setlength{\itemsep}{#2\baselineskip}}}
 {\end{list}}
\usepackage{calc}
\newcommand{\CSLBlock}[1]{\hfill\break\parbox[t]{\linewidth}{\strut\ignorespaces#1\strut}}
\newcommand{\CSLLeftMargin}[1]{\parbox[t]{\csllabelwidth}{\strut#1\strut}}
\newcommand{\CSLRightInline}[1]{\parbox[t]{\linewidth - \csllabelwidth}{\strut#1\strut}}
\newcommand{\CSLIndent}[1]{\hspace{\cslhangindent}#1}



\setlength{\emergencystretch}{3em} % prevent overfull lines

\providecommand{\tightlist}{%
  \setlength{\itemsep}{0pt}\setlength{\parskip}{0pt}}



 


\makeatletter
\@ifpackageloaded{caption}{}{\usepackage{caption}}
\AtBeginDocument{%
\ifdefined\contentsname
  \renewcommand*\contentsname{Table of contents}
\else
  \newcommand\contentsname{Table of contents}
\fi
\ifdefined\listfigurename
  \renewcommand*\listfigurename{List of Figures}
\else
  \newcommand\listfigurename{List of Figures}
\fi
\ifdefined\listtablename
  \renewcommand*\listtablename{List of Tables}
\else
  \newcommand\listtablename{List of Tables}
\fi
\ifdefined\figurename
  \renewcommand*\figurename{Figure}
\else
  \newcommand\figurename{Figure}
\fi
\ifdefined\tablename
  \renewcommand*\tablename{Table}
\else
  \newcommand\tablename{Table}
\fi
}
\@ifpackageloaded{float}{}{\usepackage{float}}
\floatstyle{ruled}
\@ifundefined{c@chapter}{\newfloat{codelisting}{h}{lop}}{\newfloat{codelisting}{h}{lop}[chapter]}
\floatname{codelisting}{Listing}
\newcommand*\listoflistings{\listof{codelisting}{List of Listings}}
\makeatother
\makeatletter
\makeatother
\makeatletter
\@ifpackageloaded{caption}{}{\usepackage{caption}}
\@ifpackageloaded{subcaption}{}{\usepackage{subcaption}}
\makeatother
\makeatletter
\@ifpackageloaded{tcolorbox}{}{\usepackage[skins,breakable]{tcolorbox}}
\makeatother
\makeatletter
\@ifundefined{shadecolor}{\definecolor{shadecolor}{HTML}{245ABE}}{}
\makeatother
\makeatletter
\@ifundefined{codebgcolor}{\definecolor{codebgcolor}{HTML}{f8f8f8}}{}
\makeatother
\makeatletter
\ifdefined\Shaded\renewenvironment{Shaded}{\begin{tcolorbox}[borderline west={3pt}{0pt}{shadecolor}, boxrule=0pt, breakable, colback={codebgcolor}, enhanced, frame hidden, sharp corners]}{\end{tcolorbox}}\fi
\makeatother

\usepackage{bookmark}
\IfFileExists{xurl.sty}{\usepackage{xurl}}{} % add URL line breaks if available
\urlstyle{same}
\hypersetup{
  pdftitle={Deploying models in production},
  pdfauthor={Giorgio Alfredo Spedicato, PhD FCAS CSPA C.Stat},
  hidelinks,
  pdfcreator={LaTeX via pandoc}}


\title{\texorpdfstring{Deploying models in
production}{Deploying models in production}}
\author{Giorgio Alfredo Spedicato, PhD FCAS CSPA C.Stat\inst{1}}
\date{2026-05-24}
\institute{Leitha SRL, Unipol Group}

\begin{document}
\frame{\titlepage}


\begin{frame}{Disclaimer}
\protect\phantomsection\label{disclaimer}
The views and opinions expressed in this presentation are those of the
author and do not necessarily reflect the position of the organization
of which he belongs.
\end{frame}

\begin{frame}{Intro}
\protect\phantomsection\label{intro}
\begin{itemize}
\tightlist
\item
  Deploying a model in production is a complex task that requires a deep
  understanding of the model, the data and the infrastructure.
\item
  Actuaries were traditionally involved in the first two aspects, but
  the third one is becoming more and more important.
\item
  Modern \emph{Python frameworks} make the third aspect \emph{easier},
  at least when presenting a POC to the stakeholders.
\end{itemize}
\end{frame}

\begin{frame}
\begin{block}{MLOps: Enhancing the ML Model Lifecycle}
\protect\phantomsection\label{mlops-enhancing-the-ml-model-lifecycle}
\begin{itemize}
\item
  \textbf{Definition}: MLOps integrates practices to automate and manage
  the entire machine learning model lifecycle, improving deployment
  efficiency and reliability in production.
\item
  \textbf{Why It's Needed}: Bridges the gap between development and
  production, enabling better collaboration between data scientists and
  engineers for models that are more performant, scalable, and
  maintainable.
\item
  \textbf{Key Benefits}: Unified workflow reduces errors, accelerates
  deployment, and provides continuous monitoring to keep models
  up-to-date and accurate.
\end{itemize}
\end{block}
\end{frame}

\begin{frame}
\begin{block}{MLOps: instruments}
\protect\phantomsection\label{mlops-instruments}
\begin{itemize}
\tightlist
\item
  \textbf{Model training}: Python pipelines, requirements files, MLflow
\item
  \textbf{Version control}: Git, GitHub, GitLab
\item
  \textbf{Model deployment}: FastAPI, Streamlit, Docker
\end{itemize}
\end{block}
\end{frame}

\begin{frame}
\begin{block}{Project presentation}
\protect\phantomsection\label{project-presentation}
\begin{itemize}
\tightlist
\item
  \textbf{Objective}: Deploying an insurance quote model in production
\item
  \textbf{Tools}:

  \begin{itemize}
  \tightlist
  \item
    \textbf{Python pipelines}: to train the frequency, severity and pure
    premium models
  \item
    \textbf{Git}: to version control the code
  \item
    \textbf{Docker}: to create a container with the models
  \item
    \textbf{FastAPI}: to deploy the models
  \item
    \textbf{Streamlit}: to create a user interface
  \end{itemize}
\end{itemize}
\end{block}
\end{frame}

\begin{frame}{General recommendations}
\protect\phantomsection\label{general-recommendations}
\begin{block}{AI tools}
\protect\phantomsection\label{ai-tools}
\begin{itemize}
\tightlist
\item
  Code assistant tools are becoming increasingly widespread in actuarial
  and data science workflows, significantly improving productivity and
  accelerating development cycles.
\item
  These tools can support developers in writing boilerplate code,
  generating documentation, debugging, and exploring alternative
  implementations.
\item
  However, AI-generated code should never be accepted blindly. Outputs
  must always be critically reviewed, validated, and tested before being
  used in production environments.
\item
  Particular attention should be paid to:

  \begin{itemize}
  \tightlist
  \item
    hidden logical errors,
  \item
    inefficient implementations,
  \item
    security vulnerabilities,
  \item
    incorrect assumptions on data structures or business rules,
  \item
    and hallucinated libraries or functions.
  \end{itemize}
\item
  In actuarial contexts, domain expertise remains essential: the final
  responsibility for the correctness, interpretability, and regulatory
  compliance of the model always lies with the practitioner.
\item
  AI tools should therefore be considered as productivity enhancers
  rather than substitutes for actuarial judgment and software
  engineering best practices.
\end{itemize}
\end{block}
\end{frame}

\begin{frame}
\begin{block}{Data exploration and understanding}
\protect\phantomsection\label{data-exploration-and-understanding}
\begin{itemize}
\tightlist
\item
  Before training or deploying any predictive model, it is fundamental
  to thoroughly explore and understand the underlying data.
\item
  Univariate and bivariate exploratory analyses are essential to:

  \begin{itemize}
  \tightlist
  \item
    understand the phenomenon being modeled,
  \item
    identify data quality issues,
  \item
    detect anomalies and outliers,
  \item
    verify variable distributions,
  \item
    assess missing values,
  \item
    and identify potential relationships between predictors and targets.
  \end{itemize}
\item
  Exploratory analysis also helps identify:

  \begin{itemize}
  \tightlist
  \item
    incorrect formats or inconsistent encodings,
  \item
    data leakage risks,
  \item
    unstable variables,
  \item
    highly imbalanced targets,
  \item
    and unexpected correlations.
  \end{itemize}
\item
  Visualization techniques such as histograms, boxplots, scatter plots,
  correlation matrices, and grouped summaries often provide insights
  that are difficult to detect through automated pipelines alone.
\item
  A good understanding of the data generation process is frequently more
  valuable than experimenting with increasingly complex models.
\end{itemize}
\end{block}
\end{frame}

\begin{frame}
\begin{block}{Production mindset}
\protect\phantomsection\label{production-mindset}
\begin{itemize}
\tightlist
\item
  A model that performs well in a notebook is not necessarily ready for
  production.
\item
  Production-ready actuarial models should emphasize:

  \begin{itemize}
  \tightlist
  \item
    reproducibility,
  \item
    monitoring,
  \item
    explainability,
  \item
    versioning,
  \item
    robustness,
  \item
    and maintainability.
  \end{itemize}
\item
  Whenever possible:

  \begin{itemize}
  \tightlist
  \item
    separate training and inference pipelines,
  \item
    track model versions and datasets,
  \item
    log predictions and input data,
  \item
    and monitor performance drift over time.
  \end{itemize}
\item
  Simpler and interpretable models are often preferable to highly
  complex solutions if they provide comparable predictive performance.
\end{itemize}
\end{block}
\end{frame}

\begin{frame}
\begin{block}{Communication and documentation}
\protect\phantomsection\label{communication-and-documentation}
\begin{itemize}
\tightlist
\item
  Clear documentation is a key component of production systems.
\item
  Assumptions, limitations, preprocessing steps, feature engineering
  choices, and validation methodologies should be explicitly documented.
\item
  Effective communication between actuaries, data scientists, software
  engineers, and business stakeholders is critical to ensure that models
  are correctly understood and properly integrated into operational
  processes.
\end{itemize}
\end{block}
\end{frame}

\begin{frame}{Python pipelines}
\protect\phantomsection\label{python-pipelines}
\begin{block}{Structure}
\protect\phantomsection\label{structure}
\begin{itemize}
\tightlist
\item
  Ingests the data, and clean
\item
  Fits the models, saves them
\item
  Assess the performance of the models, uses MLflow to log the results
  and artifacts
\end{itemize}
\end{block}
\end{frame}

\begin{frame}[fragile]
\begin{block}{Running the pipeline}
\protect\phantomsection\label{running-the-pipeline}
\begin{itemize}
\tightlist
\item
  see the \texttt{main.py} file
\item
  the directory \texttt{steps} contains the steps of the pipeline
\item
  It can be run as \texttt{python\ main.py}
\end{itemize}
\end{block}
\end{frame}

\begin{frame}[fragile]
\begin{block}{Requirements}
\protect\phantomsection\label{requirements}
\begin{itemize}
\tightlist
\item
  Python packages specified in a \texttt{requirements.txt} file
\item
  The \texttt{requirements.txt} file is used to create a virtual
  environment
\item
  The Dockerfile uses the virtual environment to create a container
\end{itemize}
\end{block}
\end{frame}

\begin{frame}[fragile]
\begin{block}{MLflow}
\protect\phantomsection\label{mlflow}
\begin{itemize}
\tightlist
\item
  MLflow (Alla and Adari 2020) is an open source platform for managing
  the end-to-end machine learning lifecycle.
\item
  It is used to log the results and artifacts of the models, that can be
  inspected in the MLFlow UI at \texttt{http://localhost:5000} after
  running the pipeline
\item
  It is organized in experiments and runs
\end{itemize}
\end{block}
\end{frame}

\begin{frame}{Streamlit walkthrough}
\protect\phantomsection\label{streamlit-walkthrough}
\begin{block}{Why streamlit?}
\protect\phantomsection\label{why-streamlit}
\begin{itemize}
\tightlist
\item
  Streamlit (Raghavendra 2022) is an open-source app framework for
  Machine Learning and Data Science projects.
\item
  It creates a user interface for the models
\item
  It is easy to use and to deploy
\end{itemize}
\end{block}
\end{frame}

\begin{frame}[fragile]
\begin{block}{Streamlit code}
\protect\phantomsection\label{streamlit-code}
\begin{itemize}
\tightlist
\item
  the \texttt{st.} functions are used to create the user interface
\item
  the session state is used to track user's choices, e.g.~button clicks
\item
  the models are loaded, cached and used to make predictions
\end{itemize}
\end{block}
\end{frame}

\begin{frame}[fragile]
\begin{block}{Running the app}
\protect\phantomsection\label{running-the-app}
\begin{itemize}
\tightlist
\item
  The app is run with \texttt{streamlit\ run\ quote-page.py}
\item
  An initial section allows to input che policyholder's data
  (pre-defined values exists)
\item
  A button click will send the data to the predict pipeline and show the
  results
\end{itemize}
\end{block}
\end{frame}

\begin{frame}[fragile]{FastAPI Walkthrough}
\protect\phantomsection\label{fastapi-walkthrough}
\begin{block}{Why FastAPI?}
\protect\phantomsection\label{why-fastapi}
\begin{itemize}
\tightlist
\item
  FastAPI (Lathkar 2023) is a modern framework for building APIs, ideal
  for deploying machine learning models.
\item
  It enables exposing models as a REST API, allowing other applications
  to use them, for example, via Python's \texttt{requests} library.
\item
  It uses Python type hints to validate the input and output of the API,
  improving code readability and maintainability.
\end{itemize}
\end{block}
\end{frame}

\begin{frame}[fragile]
\begin{block}{Structure of a FastAPI App}
\protect\phantomsection\label{structure-of-a-fastapi-app}
\begin{itemize}
\tightlist
\item
  The application is defined in a Python file, typically named
  \texttt{app.py}.
\item
  The app runs with the ASGI server \texttt{uvicorn}, allowing FastAPI
  to handle asynchronous requests, which increases scalability and
  speed.
\item
  The app structure includes a startup event for loading models,
  endpoints for handling requests, and dependencies to manage model
  usage.
\end{itemize}
\end{block}
\end{frame}

\begin{frame}[fragile]
\begin{block}{Key Components of the Insurance Prediction API}
\protect\phantomsection\label{key-components-of-the-insurance-prediction-api}
\begin{block}{1. Input Schema}
\protect\phantomsection\label{input-schema}
\begin{itemize}
\tightlist
\item
  The input is structured using the \texttt{Insured} model, a class that
  describes the characteristics of the insured person, such as vehicle
  power (\texttt{VehPower}), vehicle age (\texttt{VehAge}), driver age
  (\texttt{DrivAge}), and other attributes relevant for calculating
  insurance premiums.
\item
  Each field includes a title, description, and validity limits using
  \texttt{Field} from \texttt{pydantic}. This ensures input data
  consistency and completeness.
\end{itemize}
\end{block}
\end{block}
\end{frame}

\begin{frame}[fragile]
\begin{block}{2. Output Schema}
\protect\phantomsection\label{output-schema}
\begin{itemize}
\tightlist
\item
  The output is managed by the \texttt{PredictionResponse} class,
  defining the key variables predicted by the model: frequency
  (\texttt{Frequency}), severity (\texttt{Severity}), and pure premium
  (\texttt{Pure\_Premium}).
\item
  The \texttt{PredictionResponse} class ensures that the API always
  returns a standard format, simplifying integration with other
  applications.
\end{itemize}
\end{block}
\end{frame}

\begin{frame}[fragile]
\begin{block}{3. Startup Event}
\protect\phantomsection\label{startup-event}
\begin{itemize}
\tightlist
\item
  A FastAPI startup event, defined with
  \texttt{@app.on\_event("startup")}, loads the \texttt{CatBoost} models
  from the \texttt{models} folder. This prevents reloading the models
  every time a request is made, improving efficiency.
\item
  In this example, \texttt{model\_freq} and \texttt{model\_sev} are,
  respectively, the frequency and severity prediction models.
\end{itemize}
\end{block}
\end{frame}

\begin{frame}[fragile]
\begin{block}{4. Prediction Endpoint}
\protect\phantomsection\label{prediction-endpoint}
\begin{itemize}
\tightlist
\item
  The primary endpoint of the application is \texttt{/predict/}, which
  receives the insured person's data and returns a prediction.
\item
  Within the endpoint, the \texttt{Frequency} model and
  \texttt{Severity} model are used to calculate the pure premium:

  \begin{itemize}
  \tightlist
  \item
    \textbf{Frequency} (\texttt{Frequency}): predicts the number of
    claims.
  \item
    \textbf{Severity} (\texttt{Severity}): predicts the average cost of
    claims.
  \item
    \textbf{Pure Premium} (\texttt{Pure\_Premium}): calculated as
    \texttt{Frequency\ *\ Severity}, representing the expected premium
    for the customer.
  \end{itemize}
\item
  The endpoint uses FastAPI's \texttt{Depends} to load the model
  dependencies and then uses \texttt{Predictor} to perform predictions.
\end{itemize}
\end{block}
\end{frame}

\begin{frame}[fragile]
\begin{block}{FastAPI in Action}
\protect\phantomsection\label{fastapi-in-action}
\begin{itemize}
\tightlist
\item
  The API can be run locally using the command
  \texttt{uvicorn\ app:app\ -\/-reload}, enabling rapid development and
  testing.
\item
  Once the application is online, you can interact with the prediction
  endpoint using tools like \texttt{curl} or directly from FastAPI's
  interactive documentation available at \texttt{/docs}.
\end{itemize}
\end{block}
\end{frame}

\begin{frame}{Docker}
\protect\phantomsection\label{docker}
\begin{block}{Why Docker?}
\protect\phantomsection\label{why-docker}
\begin{itemize}
\tightlist
\item
  Docker (Jangla 2018) is a platform for developing, shipping, and
  running applications.
\item
  It allows you to create containers with models and dependencies,
  ensuring a consistent and isolated environment.
\item
  Commonly used to deploy machine learning models in production.
\end{itemize}
\end{block}
\end{frame}

\begin{frame}[fragile]
\begin{block}{Dockerfile Overview}
\protect\phantomsection\label{dockerfile-overview}
\begin{itemize}
\tightlist
\item
  The Dockerfile defines steps to create a container for our FastAPI
  application, including:

  \begin{itemize}
  \tightlist
  \item
    Loading a base Python image
  \item
    Copying models, code, and dependencies
  \item
    Installing packages from \texttt{requirements.txt}
  \item
    Starting the FastAPI app
  \end{itemize}
\end{itemize}
\end{block}
\end{frame}

\begin{frame}[fragile]
\begin{block}{Docker commands}
\protect\phantomsection\label{docker-commands}
\begin{itemize}
\tightlist
\item
  \texttt{docker\ build\ -t\ deployer\ .} to build the image
\item
  \texttt{docker\ run\ -d\ -\/-name\ deployer\ -p\ 8080:8080\ deployer:latest}
  to run the container that exposes the app on port 8000
\item
  \texttt{docker\ stop\textasciigrave{}\textasciigrave{}deployer}to stop
  the container
\end{itemize}
\end{block}
\end{frame}

\begin{frame}[allowframebreaks]{References}
\protect\phantomsection\label{references}
\protect\phantomsection\label{refs}
\begin{CSLReferences}{1}{1}
\bibitem[\citeproctext]{ref-alla2020}
Alla S, Adari SK (2020)
\href{https://doi.org/10.1007/978-1-4842-6549-9_4}{Introduction to
MLFlow}. Apress, pp 125--227

\bibitem[\citeproctext]{ref-jangla2018b}
Jangla K (2018)
\href{https://doi.org/10.1007/978-1-4842-3936-0_2}{Docker}. Apress, pp
9--17

\bibitem[\citeproctext]{ref-lathkar2023}
Lathkar M (2023)
\href{https://doi.org/10.1007/978-1-4842-9178-8_1}{Introduction to
FastAPI}. Apress, pp 1--28

\bibitem[\citeproctext]{ref-raghavendra2022}
Raghavendra S (2022)
\href{https://doi.org/10.1007/978-1-4842-8983-9_1}{Introduction to
streamlit}. Apress, pp 1--15

\end{CSLReferences}
\end{frame}




\end{document}
